\section{Introduction}
\label{sec:introduction}

Ranking-based decision making is prevalent in various application domains, including hiring~\cite{GeyikAK19} and school admission~\cite{peskun2007effectiveness}. Typically, this process involves selecting qualifying candidates based on specific criteria (e.g., for a job position) and ranking them using a quantitative measure to identify the top candidates among those who qualify (e.g., for a job interview or offer). This process may be automated and expressed using %order by 
SQL queries, with the {\tt WHERE} clause used to select candidates who meet certain requirements, and the {\tt ORDER BY} clause used to rank them. We next illustrate this idea using a simple example in the context of awarding scholarships.

% Analysts are often tasked with writing database queries whose results provide the best items from some set. For example, a company looking to fill a job opening may have time to interview only a few candidates and therefore would like to invite only the best applicants who have the desired qualifications to be interviewed. Such queries are made of two main components: the requirements that must be fulfilled by the items, and the preferences that determine which items fulfilling the requirements are more valuable than others. These components map themselves naturally to the {\tt WHERE} clause and {\tt ORDER BY} expression of SQL queries. These two SQL constructs are how many of these decision support queries are implemented in practice.

\begin{example}
	\label{ex:running}
	Consider a foundation that wishes to grant six high-performing students scholarships to universities in order to encourage participation in STEM programs. 
 \revc{The foundation utilizes a database of all students seeking scholarships provided by their schools, which may be filtered according to the requirements of the foundation.} 
 % \revc{The foundation partners with a scholarship matching service and receives a database of students seeking scholarships.} 
 % \Cref{tab:candidates} shows the students who have \revc{signed up to the service}, forming a dataset with five attributes: a unique ID, gender, family's income level, grade point average (GPA), and SAT score. 
 \Cref{tab:candidates} shows the students dataset, consisting of five attributes: a unique ID, gender, family's income level, grade point average (GPA), and SAT score. \revc{The schools also provide information on the student's 
 % \revc{Students also} submitted their 
 involvements with extracurricular activities}, which are shown in \Cref{tab:internships} as a dataset with two attributes: the student's ID and an abbreviation representing the activity in which they participated. The set of activities in \Cref{tab:internships} consists of robotics ($RB$), Science Olympiad ($SO$), Math Olympiad ($MO$), game development ($GD$), and a STEM tutoring organization ($TU$).
 % \yuval{explain what each abbreviation stands for, or don't use abbreviation}
 
    The foundation would like to award these scholarships to students who have displayed interest in STEM fields through their involvement in extracurricular activities and have maintained a minimum GPA. The selected students are ranked by their SAT exam scores, and the foundation grants funding to the best six students and additional funding to the top three students.
    % From this set of students, they are ranked by their SAT exam scores in descending order. The foundation grants funding to the best six students and grants additional funding to the top three students. 
    These requirements can be expressed using the following query, which selected students who have participated in an extracurricular robotics club with a minimum GPA of $3.7$:
    
    % In order to use the collected data, an analyst for the foundation writes the following query to retrieve students who have participated in an extracurricular robotics (abbreviated $RB$) club with a minimum GPA of 3.7:
    \begin{center}
    \footnotesize
    \begin{tabular}{l}
        \verb"SELECT DISTINCT ID, Gender, Income "\\
        \verb"FROM Students NATURAL JOIN Activities"\\
        \verb"WHERE GPA >= 3.7 AND Activity = 'RB'"\\
        \verb"ORDER BY SAT DESC"\\
    \end{tabular}
    \end{center}
    We refer to this query throughout as the \running{}. Evaluating this query over the datasets in \Cref{tab:candidates,tab:internships} produces the ranking $[t_4, t_7, t_8, t_{10}, t_{11}, t_{12}]$ %\felix{should we change these to $i$?}. \julia{I like $t_1$ etc, but I think we should change the tables to use, say, $t_1$ as the ID, not $1$. It's customary to refer to a tuple by its id (i.e., no need to list both $t_1$ and $1$, as is done in Table 4.} 
    Therefore, the foundation awards students $t_4$, $t_7$, $t_8$ with an extra scholarship and students $t_{10}$, $t_{11}$, and $t_{12}$ with the regular amount.
    % \yuval{discuss -  should we select distinct ID and not *?}
\end{example}



\begin{table}[ht]
\begin{minipage}{0.6\linewidth}
        \centering
		\caption{\candidates{}}
		\label{tab:candidates}
        \footnotesize
		\begin{tabular}{lccrr}
		\hline
		\textbf{ID} & \textbf{Gender} & \textbf{Income} & \textbf{GPA} & \textbf{SAT $\downarrow$} \\ \hline
		$t_1$           & M               & Medium         & 3.7      & 1590 \\    
		$t_2$           & F               & Low         & 3.8      & 1580 \\    
		$t_3$           & F               & Low         & 3.6      & 1570 \\    
		$t_4$           & M               & High         & 3.8      & 1560 \\    
		$t_5$           & F               & Medium         & 3.6      & 1550 \\    
		$t_6$           & F               & Low         & 3.7      & 1550 \\    
		$t_7$           & M               & Low         & 3.7      & 1540 \\    
		$t_8$           & F               & High         & 3.9      & 1530 \\    
		$t_9$           & F               & Medium         & 3.8      & 1530 \\    
		$t_{10}$           & M               & High         & 3.7      & 1520 \\    
		$t_{11}$           & F               & Low         & 3.8      & 1490 \\    
		$t_{12}$           & M               & Medium         & 4.0      & 1480 \\    
		$t_{13}$           & M               & High         & 3.5      & 1430 \\  
  	$t_{14}$           & F               & Low         & 3.7      & 1410 \\  
		\end{tabular}
\end{minipage}
\hfill
\begin{minipage}{.35\linewidth}
    \centering
    \caption{\internships{}}
    \label{tab:internships}
    \footnotesize
    \begin{tabular}{lc}
    \hline
    \textbf{ID} & \textbf{Activity} \\ \hline
        $t_{1}$            & SO        \\  
        $t_{2}$            & SO        \\  
        $t_{3}$            & GD        \\  
        $t_{4}$            & RB        \\
        $t_{4}$            & TU        \\
        $t_{5}$            & MO        \\  
        $t_{6}$            & SO        \\  
        $t_{7}$            & RB        \\  
        $t_{8}$          & RB        \\
        $t_{8}$           & TU        \\
        $t_{10}$           & RB        \\  
        $t_{11}$           & RB        \\  
        $t_{12}$           & RB        \\
        $t_{14}$          & RB        \\
    \end{tabular}
\end{minipage}
\end{table}

% \begin{table}[ht]

% \end{table}


If the query is part of some high-stake decision-making process, stating diversity requirements as cardinality constraints over the presence of some demographic groups in the top-$k$ result is natural. For instance, in the above example, the foundation may wish to promote female students in STEM by awarding a proportional number of scholarships to male and female applicants, i.e.,  top-$6$ tuples in the output should include at least three females. Moreover, to expand access to STEM education, the foundation may also wish to limit the extended scholarships granted to students from high-income families. Namely, the top-$3$ results should include at most one student with a high income. The \running{} does not satisfy these constraints since the top-$6$ tuples are $t_4, t_7, t_8, t_{10}, t_{11}$ and $t_{12}$ include only two females ($t_8$ and $t_{11}$), and the top-$3$ includes two students from high-income families ($t_4$ and $t_{8}$). 

% The outputs of these ranking queries are often used in decision-making contexts where the diversity of the results improves their usefulness. For example, in the context of \Cref{ex:running}, awarding scholarships to a more diverse group of students may assist in the foundation's mission to expand access to STEM education. Furthermore, biases have been observed in real-life deployments of such ranking processes used in decision support \cite{HBRArticle}.

% \begin{example}
%     \label{ex:lack-diversity}
%     Recall that the output of the \running{} from \Cref{ex:running} is $[t_4, t_7, t_8, t_{10}, t_{11}, t_{12}]$. In our scenario, the top three students $t_4,t_7$ and $t_8$ receive extra funding. However, students \#4 and \#8 are classified as high-income, meaning that they benefit less overall from the extra funding compared to students with medium or low-incomes. Therefore, the analyst would ideally like to limit the top three students to include at most one high income individual. In the top six students, only two students (students \#8 and \#11) are women, despite the fact women are more than half of the applicants in \Cref{tab:candidates} are women. In order to promote greater access to STEM education across all demographic groups, the foundation would like to award at least half of their scholarships to women, i.e. having at least three women in the top six best students.
% \end{example}

% Having diversity within the resulting rankings also provides sets of items that span a number of contexts or environments that would otherwise be ignored. For example, consider that a website ranking restaurants may want to refine its default set of filters (e.g. default selected neighborhoods, minimum number of reviews) for cities where those filters would be excluding restaurants that are otherwise equally well rated by reviewers (see \cite{SPDK+23} for such a scenario). As such websites are often essential in driving customers to these restaurants, diversifying the resulting ranking provides users with quality restaurants that would otherwise be excluded from exposure, benefiting both customers and proprietors.


% \subsubsection*{Ensuring diverse outputs}


% The problem of ensuring diverse outputs in ranking queries has received much recent attention from the research community~\cite{YS17, CSV18, YGS19, AJS19, KR18, CMV20}. For instance, in \cite{YS17, CSV18, YGS19}, output rankings are modified directly in a \emph{post-processing step}, in order to satisfy a given set of constraints specifying the cardinality of protected groups in the resulting ranking. E.g., to fulfill the desired constraints in the above example, the foundation may manipulate the output, awarding $t_4$, $t_5$, $t_6$, $t_7$, $t_8$ and $t_{10}$ with a scholarship, where $t_4$, $t_5$ and $t_6$ will get the extended grant.
 % However, this leaves open the question of how one may obtain such a ranking in the first place. Furthermore, approaches that modify the output ranking directly risk violating regulatory requirements by making decisions on the position of an item in the ranking (and therefore impacting its selection) based on some sensitive attributes. 
%
 In this paper, we propose a novel \emph{in-processing} method to improve the diversity of a ranking by refining the query that produces it.% to achieve an output that satisfies
 %satisfy the desired constraints.
 
 \begin{example}
    \label{ex:first_refinement}
    The \running{} may be refined by adjusting the condition on \verb"Activity" to include students involved in Science Olympiad ($SO$), resulting in the following query:
    \begin{center}
    \footnotesize
    \begin{tabular}{l}
        \verb"SELECT DISTINCT ID, Gender, Income "\\
        \verb"FROM Students NATURAL JOIN Activities"\\
        \verb"WHERE GPA >= 3.7 AND (Activity = 'RB' OR Activity = 'SO')"\\
        \verb"ORDER BY SAT DESC"\\
    \end{tabular}
    \end{center}
    Note that the essence of the query (selecting students who have displayed interest in STEM) is maintained by the refined query, while the constraints are satisfied as the top-$6$ tuples ($t_1, t_2, t_4, t_6, t_7,$ and $t_8$) consist of three women ($t_2, t_6$ and $t_8$) where the top-$3$ includes only a single student ($t_4$) with high-income.
    
    % is $[t_4, t_7, t_8, t_{10}, t_{11}, t_{12}]$. We may reformulate the \verb"WHERE" clause of the \running{} to obtain the following query: 
    
    % Evaluating this query over our dataset produces the ranking $[t_1,$ $t_2,$ $t_4,$ $t_6,$ $t_7,$ $t_8,$ $t_{10},$ $t_{11},$ $t_{12}]$. Notice now that of the top three students, only student \#4 is high-income. Furthermore, the top six students now consist of three men and three women. By adjusting the condition on \verb"Activity" in the \running{} to include students involved in Science Olympiad ($SO$), we were able to obtain a ranking that is diverse according to the requirements of the foundation.
\end{example}


 
%\julia{I modified this a bit, but also think that this paragraph can be merged with the paragraph above Example 1.2 for better flow / to get rid of redundancy.  Better to move the paragraph from before Example 1.2 here, not vice versa, to get to the example ASAP.} 

\revmulti{The notion of refining queries to satisfy a set of diversity constraints was recently presented in~\cite{ERICA, ERICAfull}, however, this work focuses on cardinality constraints over the {\it entire} output and does not consider the order of tuples}.
The problem of ensuring diverse outputs in ranking queries has received much recent attention from the research community~\cite{YS17, CSV18, YGS19, AJS19, KR18, CMV20}. For instance, in \cite{YS17, CSV18, YGS19}, output rankings are modified directly in a \emph{post-processing step}, in order to satisfy a given set of constraints over the cardinality of protected groups in the ranking. E.g., to fulfill the desired constraints in the above example, the foundation may manipulate the output, awarding $t_4$, $t_5$, $t_6$, $t_7$, $t_8$ and $t_{10}$ with a scholarship, where $t_4$, $t_5$ and $t_6$ will get the extended grant. However, this leaves open the question of \emph{how one may obtain such a ranking in the first place}.

% specifying the cardinality of protected groups in the resulting ranking. E.g., to fulfill the desired constraints in the above example, the foundation may manipulate the output, awarding $t_4$, $t_5$, $t_6$, $t_7$, $t_8$ and $t_{10}$ with a scholarship, where $t_4$, $t_5$ and $t_6$ will get the extended grant. However, this leaves open the question of how one may obtain such a ranking in the first place.

% Post-processing algorithms, on the other hand, will include or exclude individuals directly based on the group they belong to and therefore treating individuals differently.

Further, post-processing may be problematic to use to improve diversity, for two reasons: (1) by definition, it modifies the results after they were computed, raising a procedural fairness concern, and (2) it may explicitly use information about demographic or otherwise protected group membership, raising a disparate treatment concern. In contrast, in-processing %(here: modifying the query predicates) 
is usually legally permissible, essentially because it applies the same evaluation process to all individuals. \revc{That is, by modifying the query we produce a new set of requirements, and test all individuals against these same requirements.  In contrast, post-processing methods may decide to include or exclude individuals based on which groups they belong to, therefore treading individuals differently depending on group membership.}
%to which all individuals are tested equally. Post-processing methods may decide to include or exclude individuals based on which groups they belong to, violating this principle.} 
Alternative \revc{in-processing} solutions involve adjustments to the ranking algorithm~\cite{AJS19} or modifying items to produce a different score~\cite{KR18, CMV20}. Our approach, conversely, assumes the ranking algorithms and scores of different items are well-designed and fixed, and we aim to modify the set of tuples to be ranked.
 % Previous work has addressed the concerns we have mentioned so far by instead making modifications to specific classes of ranking algorithms as in \cite{AJS19} or modifying the items themselves to produce a different score as in \cite{KR18, CMV20}. These approaches are known more generally as {\it pre-processing} and {\it in-processing} approaches, as described in \cite{ZYS23}. However, existing approaches assume that the entire input must be ranked. In our proposed framework, we refine the set of requirements made on the items, therefore modifying which items are input to the ranking algorithm, instead of adjusting the ranking algorithm or the items themselves. 

 Our goal is to find minimal refinements to the original query that fulfill a specified set of constraints, however, we note that the notion of minimality may be defined in different ways, depending, for example, on the legal requirements or on the user's preferences.

 \begin{example}
    \label{ex:refine-and-distance}
%    Recall that the output of the \running{} from \Cref{ex:running} is $[t_4, t_7, t_8, t_{10}, t_{11}, t_{12}]$. We may reformulate the \verb"WHERE" clause of the \running{} to obtain the following query:
    % \begin{center}
    % \begin{small}
    % \begin{tabular}{l}
    %     \verb"SELECT * "\\
    %     \verb"FROM Students NATURAL JOIN Activities"\\
    %     \verb"WHERE Gpa >= 3.7 AND (Activity = 'RB' OR Activity = 'SO')"\\
    %     \verb"ORDER BY Sat DESC"\\
    % \end{tabular}
    % \end{small}
    % \end{center}
    % Evaluating this query over our dataset produces the ranking $[t_1,$ $t_2,$ $t_4,$ $t_6,$ $t_7,$ $t_8,$ $t_{10},$ $t_{11},$ $t_{12}]$. Notice now that of the top three students, only student \#4 is high-income. Furthermore, the top six students now consist of three men and three women. By adjusting the condition on \verb"Activity" in the \running{} to include students involved in Science Olympiad ($SO$), we were able to obtain a ranking that is diverse according to the requirements of the foundation.
    % Similar to \Cref{ex:first_refinement}, 
    We may refine the \running{} by relaxing the GPA requirement to $3.6$ and including students who participated in a game development activity $(GD)$,  obtaining the following query: 
    % We obtain the following query by making such refinements:
    \begin{center}
    \footnotesize
    \begin{tabular}{l}
        \verb"SELECT DISTINCT ID, Gender, Income "\\
        \verb"FROM Students NATURAL JOIN Activities"\\
        \verb"WHERE GPA >= 3.6 AND (Activity = 'RB' OR Activity = 'GD')"\\
        \verb"ORDER BY SAT DESC"\\
    \end{tabular}
    \end{center}
    Similarly to the refined query from Example~\ref{ex:first_refinement}, the top-$6$ students % that are selected 
    ($t_3$, $t_4$, $t_7$, $t_8$, $t_{10}$, $t_{11}$, and $t_{12}$) include three women ($t_3$, $t_8$, and $t_{11}$), and there is only a single high-income student ($t_3$) among the top-$3$. While the predicates of this refined query are intuitively more distant from the original query than our prior refinement in Example~\ref{ex:first_refinement} (two modifications compared to a single one), its output is more similar to the output of the original query (the top-$3$ sets differ by one tuple).
    
    % However, the predicates of this refined query are intuitively more distant from the original query than our prior refinement. Specifically, even though we add one activity to both queries, we also modify the GPA predicate in this refinement. On the other hand, the set of the top three of students of this refinement is intuitively more similar to the set of the top three students of the original \running{}. While this refinement has two of the original students (students \#4 and \#7) and only one new student (student \#3) in the top-3, the prior refinement has only one new student (student \#4). Therefore, even though both of these refinements are diverse enough according to the foundation's requirements, they are closer than one another to the original \running{} in their own senses. 
\end{example}

To accommodate different alternative query relaxation objectives, as illustrated above, we propose a framework that allows the user to specify their preferred notion of minimality. 

To the best of our knowledge, our work is the first to intervene on the ranking process by modifying {\it which} items are being considered by the ranking algorithm. This admits a large class of ranking algorithms while keeping the relative order of tuples consistent. This---of course---does not come for free; the coarseness of refinements means that there may be no refinement that produces a satisfactorily diverse ranking. Therefore, we study the problem of finding a refined query %with a refined set of requirements 
that is within a specified maximum distance from satisfying all of the constraints, if one exists. %, while maintaining the intent of the query as much as possible. \julia{this is redundant, streamlined the sentence}


% We present the problem of refining the selection predicates of a SPUJ order-by query to meet constraints regarding the representation of various groups in the top-$k$ positions of the query's output (Section~\ref{sec:background}). Our goal is to identify the smallest refinements needed to satisfy these constraints, acknowledging that there are different ways to define minimality in this context.






\paragraph*{Contributions \& roadmap}
We begin by formalizing the \problem{} problem of obtaining a refined query that is closest, according to a given distance measure, to the original query, while still adhering to a set of cardinality constraints within a maximum distance, and show that this problem is {\sf NP-hard} (\Cref{sec:background}).
We thus propose modeling the problem as a Mixed Integer Linear Program (MILP) and utilizing it to derive an approximate solution (\Cref{sec:search}). Inspired by the use of data annotations (provenance) to perform what-if analysis, i.e., to efficiently reevaluate queries using algebraic expression without constantly accessing a DBMS (see, e.g., ~\cite{DIMT13, BourhisDM16, DeutchMT14, MoskovitchLJ22}), we construct the MILP using data annotation variables. 
This formulation offers two advantages: it allows us to leverage the effectiveness of existing MILP solvers while avoiding the costly reevaluation of queries on the DBMS.

Existing MILP solvers may solve the problem efficiently however their performance is sensitive to the size of the program. The program generated in Section~\ref{sec:search} is linear in the data size, which can be challenging in real-life scenarios as we demonstrate in our experimental evaluation. We, therefore, propose optimizations to make our approach more scalable, using the relevancy of the data and the structure of the set of cardinality constraints to prune and relax our problem (\Cref{sec:optimizations}). In \Cref{sec:experiments}, we present an extensive experimental evaluation. We developed a dedicated benchmark consisting of real-life datasets and considering realistic scenarios. Our results show the efficiency and scalability of our approach with respect to different parameters of the problem. 
% \revmulti{Finally, we compare and contrast our framework with Erica \cite{MLJ22,ERICA,ERICAfull}, a system that finds refinements for cardinality constraints on the {\it entire} output instead of the top-$k$. We show that our solution can provide closer refinements with better performance for our instances of our problem that overlap with their system.}





% (in some user defined measure, as motivated by \Cref{ex:refine-and-distance}) while remaining within some maximum distance from the set of cardinality constraints used to intervene on the output ranking in \Cref{sec:background}. Furthermore, we show that determining whether such a refined query exists is {\sf NP-complete} and therefore determine the hardness of the problem of finding the best approximation refinement. \\
% \indent In \Cref{sec:search}, we discuss the framework we have developed in order to solve the problem we study. We use ideas from data provenance to annotate tuples from the database, and leverage these annotated tuples with techniques from what-if analysis (like \cite{DIMT13}) in order to efficiently evaluate refinements without repeatedly evaluating queries across the DBMS. We also propose two optimizations to make our approach more scalable, using the relevancy of the data and the structure of the set of cardinality constraints to prune and relax our problem. \\
% \indent Finally, we present an investigation into the effectiveness and efficiency of our approach in \Cref{sec:experiments}. We use real world and well-known benchmark datasets to determine that our method is a practical intervention for ensuring diversity in rankings. \felix{conjecture, also place experiment results/conclusions} 
%We now sketch the remainder of the paper. In \Cref{sec:background}, we outline our approach and then formally define it, providing the necessary background and terminology as well. \Cref{sec:search} describes how we use this model to find the best approximation refinement. We also provide an experimental investigation into the scalability and effectiveness of our approach in \Cref{sec:experiments}. Finally, \Cref{sec:related} surveys work related to ours. 
Ranking-based decision making is prevalent in various application domains, including hiring~\cite{GeyikAK19} and school admission~\cite{peskun2007effectiveness}. Typically, this process involves selecting qualifying candidates based on specific criteria (e.g., for a job position) and ranking them using a quantitative measure to identify the top candidates among those who qualify (e.g., for a job interview or offer). This process may be automated and expressed using %order by 
SQL queries, with the {\tt WHERE} clause used to select candidates who meet certain requirements, and the {\tt ORDER BY} clause used to rank them. We next illustrate this idea using a simple example in the context of awarding scholarships.

% Analysts are often tasked with writing database queries whose results provide the best items from some set. For example, a company looking to fill a job opening may have time to interview only a few candidates and therefore would like to invite only the best applicants who have the desired qualifications to be interviewed. Such queries are made of two main components: the requirements that must be fulfilled by the items, and the preferences that determine which items fulfilling the requirements are more valuable than others. These components map themselves naturally to the {\tt WHERE} clause and {\tt ORDER BY} expression of SQL queries. These two SQL constructs are how many of these decision support queries are implemented in practice.

\begin{example}
	\label{ex:running}
	Consider a foundation that wishes to grant six high-performing students scholarships to universities in order to encourage participation in STEM programs. 
 \revc{The foundation utilizes a database of all students seeking scholarships provided by their schools, which may be filtered according to the requirements of the foundation.} 
 % \revc{The foundation partners with a scholarship matching service and receives a database of students seeking scholarships.} 
 % \Cref{tab:candidates} shows the students who have \revc{signed up to the service}, forming a dataset with five attributes: a unique ID, gender, family's income level, grade point average (GPA), and SAT score. 
 \Cref{tab:candidates} shows the students dataset, consisting of five attributes: a unique ID, gender, family's income level, grade point average (GPA), and SAT score. \revc{The schools also provide information on the student's 
 % \revc{Students also} submitted their 
 involvements with extracurricular activities}, which are shown in \Cref{tab:internships} as a dataset with two attributes: the student's ID and an abbreviation representing the activity in which they participated. The set of activities in \Cref{tab:internships} consists of robotics ($RB$), Science Olympiad ($SO$), Math Olympiad ($MO$), game development ($GD$), and a STEM tutoring organization ($TU$).
 % \yuval{explain what each abbreviation stands for, or don't use abbreviation}
 
    The foundation would like to award these scholarships to students who have displayed interest in STEM fields through their involvement in extracurricular activities and have maintained a minimum GPA. The selected students are ranked by their SAT exam scores, and the foundation grants funding to the best six students and additional funding to the top three students.
    % From this set of students, they are ranked by their SAT exam scores in descending order. The foundation grants funding to the best six students and grants additional funding to the top three students. 
    These requirements can be expressed using the following query, which selected students who have participated in an extracurricular robotics club with a minimum GPA of $3.7$:
    
    % In order to use the collected data, an analyst for the foundation writes the following query to retrieve students who have participated in an extracurricular robotics (abbreviated $RB$) club with a minimum GPA of 3.7:
    \begin{center}
    \footnotesize
    \begin{tabular}{l}
        \verb"SELECT DISTINCT ID, Gender, Income "\\
        \verb"FROM Students NATURAL JOIN Activities"\\
        \verb"WHERE GPA >= 3.7 AND Activity = 'RB'"\\
        \verb"ORDER BY SAT DESC"\\
    \end{tabular}
    \end{center}
    We refer to this query throughout as the \running{}. Evaluating this query over the datasets in \Cref{tab:candidates,tab:internships} produces the ranking $[t_4, t_7, t_8, t_{10}, t_{11}, t_{12}]$ %\felix{should we change these to $i$?}. \julia{I like $t_1$ etc, but I think we should change the tables to use, say, $t_1$ as the ID, not $1$. It's customary to refer to a tuple by its id (i.e., no need to list both $t_1$ and $1$, as is done in Table 4.} 
    Therefore, the foundation awards students $t_4$, $t_7$, $t_8$ with an extra scholarship and students $t_{10}$, $t_{11}$, and $t_{12}$ with the regular amount.
    % \yuval{discuss -  should we select distinct ID and not *?}
\end{example}



\begin{table}[ht]
\begin{minipage}{0.6\linewidth}
        \centering
		\caption{\candidates{}}
		\label{tab:candidates}
        \footnotesize
		\begin{tabular}{lccrr}
		\hline
		\textbf{ID} & \textbf{Gender} & \textbf{Income} & \textbf{GPA} & \textbf{SAT $\downarrow$} \\ \hline
		$t_1$           & M               & Medium         & 3.7      & 1590 \\    
		$t_2$           & F               & Low         & 3.8      & 1580 \\    
		$t_3$           & F               & Low         & 3.6      & 1570 \\    
		$t_4$           & M               & High         & 3.8      & 1560 \\    
		$t_5$           & F               & Medium         & 3.6      & 1550 \\    
		$t_6$           & F               & Low         & 3.7      & 1550 \\    
		$t_7$           & M               & Low         & 3.7      & 1540 \\    
		$t_8$           & F               & High         & 3.9      & 1530 \\    
		$t_9$           & F               & Medium         & 3.8      & 1530 \\    
		$t_{10}$           & M               & High         & 3.7      & 1520 \\    
		$t_{11}$           & F               & Low         & 3.8      & 1490 \\    
		$t_{12}$           & M               & Medium         & 4.0      & 1480 \\    
		$t_{13}$           & M               & High         & 3.5      & 1430 \\  
  	$t_{14}$           & F               & Low         & 3.7      & 1410 \\  
		\end{tabular}
\end{minipage}
\hfill
\begin{minipage}{.35\linewidth}
    \centering
    \caption{\internships{}}
    \label{tab:internships}
    \footnotesize
    \begin{tabular}{lc}
    \hline
    \textbf{ID} & \textbf{Activity} \\ \hline
        $t_{1}$            & SO        \\  
        $t_{2}$            & SO        \\  
        $t_{3}$            & GD        \\  
        $t_{4}$            & RB        \\
        $t_{4}$            & TU        \\
        $t_{5}$            & MO        \\  
        $t_{6}$            & SO        \\  
        $t_{7}$            & RB        \\  
        $t_{8}$          & RB        \\
        $t_{8}$           & TU        \\
        $t_{10}$           & RB        \\  
        $t_{11}$           & RB        \\  
        $t_{12}$           & RB        \\
        $t_{14}$          & RB        \\
    \end{tabular}
\end{minipage}
\end{table}

% \begin{table}[ht]

% \end{table}


If the query is part of some high-stake decision-making process, stating diversity requirements as cardinality constraints over the presence of some demographic groups in the top-$k$ result is natural. For instance, in the above example, the foundation may wish to promote female students in STEM by awarding a proportional number of scholarships to male and female applicants, i.e.,  top-$6$ tuples in the output should include at least three females. Moreover, to expand access to STEM education, the foundation may also wish to limit the extended scholarships granted to students from high-income families. Namely, the top-$3$ results should include at most one student with a high income. The \running{} does not satisfy these constraints since the top-$6$ tuples are $t_4, t_7, t_8, t_{10}, t_{11}$ and $t_{12}$ include only two females ($t_8$ and $t_{11}$), and the top-$3$ includes two students from high-income families ($t_4$ and $t_{8}$). 

% The outputs of these ranking queries are often used in decision-making contexts where the diversity of the results improves their usefulness. For example, in the context of \Cref{ex:running}, awarding scholarships to a more diverse group of students may assist in the foundation's mission to expand access to STEM education. Furthermore, biases have been observed in real-life deployments of such ranking processes used in decision support \cite{HBRArticle}.

% \begin{example}
%     \label{ex:lack-diversity}
%     Recall that the output of the \running{} from \Cref{ex:running} is $[t_4, t_7, t_8, t_{10}, t_{11}, t_{12}]$. In our scenario, the top three students $t_4,t_7$ and $t_8$ receive extra funding. However, students \#4 and \#8 are classified as high-income, meaning that they benefit less overall from the extra funding compared to students with medium or low-incomes. Therefore, the analyst would ideally like to limit the top three students to include at most one high income individual. In the top six students, only two students (students \#8 and \#11) are women, despite the fact women are more than half of the applicants in \Cref{tab:candidates} are women. In order to promote greater access to STEM education across all demographic groups, the foundation would like to award at least half of their scholarships to women, i.e. having at least three women in the top six best students.
% \end{example}

% Having diversity within the resulting rankings also provides sets of items that span a number of contexts or environments that would otherwise be ignored. For example, consider that a website ranking restaurants may want to refine its default set of filters (e.g. default selected neighborhoods, minimum number of reviews) for cities where those filters would be excluding restaurants that are otherwise equally well rated by reviewers (see \cite{SPDK+23} for such a scenario). As such websites are often essential in driving customers to these restaurants, diversifying the resulting ranking provides users with quality restaurants that would otherwise be excluded from exposure, benefiting both customers and proprietors.


% \subsubsection*{Ensuring diverse outputs}


% The problem of ensuring diverse outputs in ranking queries has received much recent attention from the research community~\cite{YS17, CSV18, YGS19, AJS19, KR18, CMV20}. For instance, in \cite{YS17, CSV18, YGS19}, output rankings are modified directly in a \emph{post-processing step}, in order to satisfy a given set of constraints specifying the cardinality of protected groups in the resulting ranking. E.g., to fulfill the desired constraints in the above example, the foundation may manipulate the output, awarding $t_4$, $t_5$, $t_6$, $t_7$, $t_8$ and $t_{10}$ with a scholarship, where $t_4$, $t_5$ and $t_6$ will get the extended grant.
 % However, this leaves open the question of how one may obtain such a ranking in the first place. Furthermore, approaches that modify the output ranking directly risk violating regulatory requirements by making decisions on the position of an item in the ranking (and therefore impacting its selection) based on some sensitive attributes. 
%
 In this paper, we propose a novel \emph{in-processing} method to improve the diversity of a ranking by refining the query that produces it.% to achieve an output that satisfies
 %satisfy the desired constraints.
 
 \begin{example}
    \label{ex:first_refinement}
    The \running{} may be refined by adjusting the condition on \verb"Activity" to include students involved in Science Olympiad ($SO$), resulting in the following query:
    \begin{center}
    \footnotesize
    \begin{tabular}{l}
        \verb"SELECT DISTINCT ID, Gender, Income "\\
        \verb"FROM Students NATURAL JOIN Activities"\\
        \verb"WHERE GPA >= 3.7 AND (Activity = 'RB' OR Activity = 'SO')"\\
        \verb"ORDER BY SAT DESC"\\
    \end{tabular}
    \end{center}
    Note that the essence of the query (selecting students who have displayed interest in STEM) is maintained by the refined query, while the constraints are satisfied as the top-$6$ tuples ($t_1, t_2, t_4, t_6, t_7,$ and $t_8$) consist of three women ($t_2, t_6$ and $t_8$) where the top-$3$ includes only a single student ($t_4$) with high-income.
    
    % is $[t_4, t_7, t_8, t_{10}, t_{11}, t_{12}]$. We may reformulate the \verb"WHERE" clause of the \running{} to obtain the following query: 
    
    % Evaluating this query over our dataset produces the ranking $[t_1,$ $t_2,$ $t_4,$ $t_6,$ $t_7,$ $t_8,$ $t_{10},$ $t_{11},$ $t_{12}]$. Notice now that of the top three students, only student \#4 is high-income. Furthermore, the top six students now consist of three men and three women. By adjusting the condition on \verb"Activity" in the \running{} to include students involved in Science Olympiad ($SO$), we were able to obtain a ranking that is diverse according to the requirements of the foundation.
\end{example}


 
%\julia{I modified this a bit, but also think that this paragraph can be merged with the paragraph above Example 1.2 for better flow / to get rid of redundancy.  Better to move the paragraph from before Example 1.2 here, not vice versa, to get to the example ASAP.} 

\revmulti{The notion of refining queries to satisfy a set of diversity constraints was recently presented in~\cite{ERICA, ERICAfull}, however, this work focuses on cardinality constraints over the {\it entire} output and does not consider the order of tuples}.
The problem of ensuring diverse outputs in ranking queries has received much recent attention from the research community~\cite{YS17, CSV18, YGS19, AJS19, KR18, CMV20}. For instance, in \cite{YS17, CSV18, YGS19}, output rankings are modified directly in a \emph{post-processing step}, in order to satisfy a given set of constraints over the cardinality of protected groups in the ranking. E.g., to fulfill the desired constraints in the above example, the foundation may manipulate the output, awarding $t_4$, $t_5$, $t_6$, $t_7$, $t_8$ and $t_{10}$ with a scholarship, where $t_4$, $t_5$ and $t_6$ will get the extended grant. However, this leaves open the question of \emph{how one may obtain such a ranking in the first place}.

% specifying the cardinality of protected groups in the resulting ranking. E.g., to fulfill the desired constraints in the above example, the foundation may manipulate the output, awarding $t_4$, $t_5$, $t_6$, $t_7$, $t_8$ and $t_{10}$ with a scholarship, where $t_4$, $t_5$ and $t_6$ will get the extended grant. However, this leaves open the question of how one may obtain such a ranking in the first place.

% Post-processing algorithms, on the other hand, will include or exclude individuals directly based on the group they belong to and therefore treating individuals differently.

Further, post-processing may be problematic to use to improve diversity, for two reasons: (1) by definition, it modifies the results after they were computed, raising a procedural fairness concern, and (2) it may explicitly use information about demographic or otherwise protected group membership, raising a disparate treatment concern. In contrast, in-processing %(here: modifying the query predicates) 
is usually legally permissible, essentially because it applies the same evaluation process to all individuals. \revc{That is, by modifying the query we produce a new set of requirements, and test all individuals against these same requirements.  In contrast, post-processing methods may decide to include or exclude individuals based on which groups they belong to, therefore treading individuals differently depending on group membership.}
%to which all individuals are tested equally. Post-processing methods may decide to include or exclude individuals based on which groups they belong to, violating this principle.} 
Alternative \revc{in-processing} solutions involve adjustments to the ranking algorithm~\cite{AJS19} or modifying items to produce a different score~\cite{KR18, CMV20}. Our approach, conversely, assumes the ranking algorithms and scores of different items are well-designed and fixed, and we aim to modify the set of tuples to be ranked.
 % Previous work has addressed the concerns we have mentioned so far by instead making modifications to specific classes of ranking algorithms as in \cite{AJS19} or modifying the items themselves to produce a different score as in \cite{KR18, CMV20}. These approaches are known more generally as {\it pre-processing} and {\it in-processing} approaches, as described in \cite{ZYS23}. However, existing approaches assume that the entire input must be ranked. In our proposed framework, we refine the set of requirements made on the items, therefore modifying which items are input to the ranking algorithm, instead of adjusting the ranking algorithm or the items themselves. 

 Our goal is to find minimal refinements to the original query that fulfill a specified set of constraints, however, we note that the notion of minimality may be defined in different ways, depending, for example, on the legal requirements or on the user's preferences.

 \begin{example}
    \label{ex:refine-and-distance}
%    Recall that the output of the \running{} from \Cref{ex:running} is $[t_4, t_7, t_8, t_{10}, t_{11}, t_{12}]$. We may reformulate the \verb"WHERE" clause of the \running{} to obtain the following query:
    % \begin{center}
    % \begin{small}
    % \begin{tabular}{l}
    %     \verb"SELECT * "\\
    %     \verb"FROM Students NATURAL JOIN Activities"\\
    %     \verb"WHERE Gpa >= 3.7 AND (Activity = 'RB' OR Activity = 'SO')"\\
    %     \verb"ORDER BY Sat DESC"\\
    % \end{tabular}
    % \end{small}
    % \end{center}
    % Evaluating this query over our dataset produces the ranking $[t_1,$ $t_2,$ $t_4,$ $t_6,$ $t_7,$ $t_8,$ $t_{10},$ $t_{11},$ $t_{12}]$. Notice now that of the top three students, only student \#4 is high-income. Furthermore, the top six students now consist of three men and three women. By adjusting the condition on \verb"Activity" in the \running{} to include students involved in Science Olympiad ($SO$), we were able to obtain a ranking that is diverse according to the requirements of the foundation.
    % Similar to \Cref{ex:first_refinement}, 
    We may refine the \running{} by relaxing the GPA requirement to $3.6$ and including students who participated in a game development activity $(GD)$,  obtaining the following query: 
    % We obtain the following query by making such refinements:
    \begin{center}
    \footnotesize
    \begin{tabular}{l}
        \verb"SELECT DISTINCT ID, Gender, Income "\\
        \verb"FROM Students NATURAL JOIN Activities"\\
        \verb"WHERE GPA >= 3.6 AND (Activity = 'RB' OR Activity = 'GD')"\\
        \verb"ORDER BY SAT DESC"\\
    \end{tabular}
    \end{center}
    Similarly to the refined query from Example~\ref{ex:first_refinement}, the top-$6$ students % that are selected 
    ($t_3$, $t_4$, $t_7$, $t_8$, $t_{10}$, $t_{11}$, and $t_{12}$) include three women ($t_3$, $t_8$, and $t_{11}$), and there is only a single high-income student ($t_3$) among the top-$3$. While the predicates of this refined query are intuitively more distant from the original query than our prior refinement in Example~\ref{ex:first_refinement} (two modifications compared to a single one), its output is more similar to the output of the original query (the top-$3$ sets differ by one tuple).
    
    % However, the predicates of this refined query are intuitively more distant from the original query than our prior refinement. Specifically, even though we add one activity to both queries, we also modify the GPA predicate in this refinement. On the other hand, the set of the top three of students of this refinement is intuitively more similar to the set of the top three students of the original \running{}. While this refinement has two of the original students (students \#4 and \#7) and only one new student (student \#3) in the top-3, the prior refinement has only one new student (student \#4). Therefore, even though both of these refinements are diverse enough according to the foundation's requirements, they are closer than one another to the original \running{} in their own senses. 
\end{example}

To accommodate different alternative query relaxation objectives, as illustrated above, we propose a framework that allows the user to specify their preferred notion of minimality. 

To the best of our knowledge, our work is the first to intervene on the ranking process by modifying {\it which} items are being considered by the ranking algorithm. This admits a large class of ranking algorithms while keeping the relative order of tuples consistent. This---of course---does not come for free; the coarseness of refinements means that there may be no refinement that produces a satisfactorily diverse ranking. Therefore, we study the problem of finding a refined query %with a refined set of requirements 
that is within a specified maximum distance from satisfying all of the constraints, if one exists. %, while maintaining the intent of the query as much as possible. \julia{this is redundant, streamlined the sentence}


% We present the problem of refining the selection predicates of a SPUJ order-by query to meet constraints regarding the representation of various groups in the top-$k$ positions of the query's output (Section~\ref{sec:background}). Our goal is to identify the smallest refinements needed to satisfy these constraints, acknowledging that there are different ways to define minimality in this context.






\paragraph*{Contributions \& roadmap}
We begin by formalizing the \problem{} problem of obtaining a refined query that is closest, according to a given distance measure, to the original query, while still adhering to a set of cardinality constraints within a maximum distance, and show that this problem is {\sf NP-hard} (\Cref{sec:background}).
We thus propose modeling the problem as a Mixed Integer Linear Program (MILP) and utilizing it to derive an approximate solution (\Cref{sec:search}). Inspired by the use of data annotations (provenance) to perform what-if analysis, i.e., to efficiently reevaluate queries using algebraic expression without constantly accessing a DBMS (see, e.g., ~\cite{DIMT13, BourhisDM16, DeutchMT14, MoskovitchLJ22}), we construct the MILP using data annotation variables. 
This formulation offers two advantages: it allows us to leverage the effectiveness of existing MILP solvers while avoiding the costly reevaluation of queries on the DBMS.

Existing MILP solvers may solve the problem efficiently however their performance is sensitive to the size of the program. The program generated in Section~\ref{sec:search} is linear in the data size, which can be challenging in real-life scenarios as we demonstrate in our experimental evaluation. We, therefore, propose optimizations to make our approach more scalable, using the relevancy of the data and the structure of the set of cardinality constraints to prune and relax our problem (\Cref{sec:optimizations}). In \Cref{sec:experiments}, we present an extensive experimental evaluation. We developed a dedicated benchmark consisting of real-life datasets and considering realistic scenarios. Our results show the efficiency and scalability of our approach with respect to different parameters of the problem. 
% \revmulti{Finally, we compare and contrast our framework with Erica \cite{MLJ22,ERICA,ERICAfull}, a system that finds refinements for cardinality constraints on the {\it entire} output instead of the top-$k$. We show that our solution can provide closer refinements with better performance for our instances of our problem that overlap with their system.}





% (in some user defined measure, as motivated by \Cref{ex:refine-and-distance}) while remaining within some maximum distance from the set of cardinality constraints used to intervene on the output ranking in \Cref{sec:background}. Furthermore, we show that determining whether such a refined query exists is {\sf NP-complete} and therefore determine the hardness of the problem of finding the best approximation refinement. \\
% \indent In \Cref{sec:search}, we discuss the framework we have developed in order to solve the problem we study. We use ideas from data provenance to annotate tuples from the database, and leverage these annotated tuples with techniques from what-if analysis (like \cite{DIMT13}) in order to efficiently evaluate refinements without repeatedly evaluating queries across the DBMS. We also propose two optimizations to make our approach more scalable, using the relevancy of the data and the structure of the set of cardinality constraints to prune and relax our problem. \\
% \indent Finally, we present an investigation into the effectiveness and efficiency of our approach in \Cref{sec:experiments}. We use real world and well-known benchmark datasets to determine that our method is a practical intervention for ensuring diversity in rankings. \felix{conjecture, also place experiment results/conclusions} 
%We now sketch the remainder of the paper. In \Cref{sec:background}, we outline our approach and then formally define it, providing the necessary background and terminology as well. \Cref{sec:search} describes how we use this model to find the best approximation refinement. We also provide an experimental investigation into the scalability and effectiveness of our approach in \Cref{sec:experiments}. Finally, \Cref{sec:related} surveys work related to ours. 